\section{Approximation de $\pi$}

On prend la suite

$$
\begin{cases}
u(0) &= 1 \\
v(0) &= 2 \\
u(k+1) &= \frac{u(k)v(k)}{2} \\
u(k+1) &= \sqrt{u(k+1) v(k)} \\
\end{cases}
$$

On admet que $v(k) \xrightarrow{k \to \infty} \frac{\sqrt{27}}{\pi}$.

\subsection*{Question 1}
 Ecrire un programme renvoyant pour un entier $N$ l'appoximation du nombre $\pi$ obtenue à partir de la valeur de $v(N)$.
 
\subsection*{Question 2}

Ecrire uun programme renvoyant pour uin réel $\epsilon$ strictement positif l'approximation de $\pi$ obtenue à partir de la valeur $v(n_0)$, où $n_0$ est le premier entier tel que:

$$
|\frac{u(n_0) - v(n_0)}{u(n_0) + v(n_0)}|< \epsilon
$$
